\chapter{Slicing and the Plateau direct method}\label{sec:current-slicing}
Restricting a current to a sublevel set creates a new boundary on the level set. Slicing identifies this boundary for almost every level and estimates its mass. The statements below use real-valued Lipschitz slicing functions and integer coefficients, following \cite{Sim18}, Chapter 6, Section 4, pp.~180--185.

\begin{definition}[Slice by a scalar function]\label{def:current-slice}
Let $T$ be a locally normal $m$-current, $m\geq1$, and let $f:U\to\mathbb R$ be locally Lipschitz. Its lower slice, initially defined as a distributional current, is
\[
 \langle T,f,t\rangle
 =\partial(T\llcorner\{f<t\})-(\partial T)\llcorner\{f<t\}.
\]
This formula alone does not give finite mass at every level. The almost-everywhere mass and rectifiability statements are the content of the slicing theorem. The boundary sign is fixed by restricting to $\{f<t\}$; compare \cite{Sim18}, equations (4.6)--(4.7).
\end{definition}

\begin{theorem}[Slicing of integral currents]\label{thm:integral-current-slicing}
Let $T=\llbracket M,\theta,\tau\rrbracket$ be a locally integral $m$-current in $U$, $m\geq1$, and let $f:U\to\mathbb R$ be locally Lipschitz. For almost every $t$, the slice $\langle T,f,t\rangle$ is a locally integral $(m-1)$-current supported in $\operatorname{spt}T\cap f^{-1}(t)$; the restriction $T\llcorner\{f<t\}$ is locally integral as well. For a Borel set $E\subset U$,
\[
 \int_{\mathbb R}\|\langle T,f,t\rangle\|(E)\,dt
 =\int_E|\nabla^M f|\,d\|T\|.
\]
In particular, if $f$ is $L$-Lipschitz, this is at most $L\|T\|(E)$. On the rectifiable part of $M\cap f^{-1}(t)$ where $\nabla^M f\ne0$, the slice carries multiplicity $\theta$ and orientation $\tau\mathbin{\lrcorner}(\nabla^M f/|\nabla^M f|)$, where
$\langle\tau\mathbin{\lrcorner}v,\eta\rangle=\langle\tau,v\wedge\eta\rangle$.

Sources: \cite{Sim18}, Definition 4.4 and Lemma 4.5, pp.~181--183. The proof below identifies the geometric slice with the distributional boundary in Definition \ref{def:current-slice}, including a single exceptional null set for all test forms. Rectifiable coarea and the almost-everywhere fiber tangent theorem are now proved in Theorem \ref{thm:rectifiable-coarea-chart-proof}. The manifold decomposition is proved in Theorem \ref{thm:rectifiable-c1-decomposition}, with Whitney extension and the Euclidean area/coarea formulas retained as inputs. All conclusions concerning a slice are almost-everywhere in the level.
\end{theorem}
\begin{proof}
Proposition \ref{prop:geometric-slice-integration} constructs the geometric slices, fixes their orientations and proves the mass identity. Its dimension-one endpoint uses the separately proved curve coarea statement. Proposition \ref{prop:slice-common-exceptional-set} identifies these currents with the distributional slices outside one null set, using the Lipschitz scalar product rule and countable $C^1$ tests. Proposition \ref{prop:integral-geometric-cut} then proves integrality of both the restrictions and slices, including the boundary sign.

The support assertion follows locally from the distributional identity. A test supported in $\{f<t\}$ sees the cancellation $\partial T-\partial T$; a test supported in $\{f>t\}$ sees zero in both terms. A partition of unity extends this to tests supported off the level, and tests supported off $\operatorname{spt}T$ vanish as well. Finally $|\nabla^Mf|\leq L$ gives the Lipschitz mass bound. None of these steps assumes finite global mass when only local integrality is given.
\end{proof}

\begin{proposition}[Boundary and slicing have opposite signs]\label{prop:slice-boundary-sign}
If $m\geq2$, then at almost every level for a locally integral current,
\[
 \partial\langle T,f,t\rangle=-\langle\partial T,f,t\rangle.
\]
If $m=1$, the slice is a zero-current and its boundary is zero by convention. See \cite{Sim18}, Lemma 4.5(3).
\end{proposition}
\begin{proof}
Apply $\partial$ to the definition of the slice. The boundary-squared term vanishes, leaving $-\partial((\partial T)\llcorner\{f<t\})$. Apply the same definition to $\partial T$ and use $\partial^2T=0$ to identify it with the negative slice of $\partial T$. The almost-everywhere qualification ensures the finite-mass and integral conclusions for both slices.
\end{proof}

\begin{proposition}[Choosing a radius with controlled cutting boundary]\label{prop:current-good-radius}
Let $T$ be locally integral near $\overline B_R(a)$ and let $0<r<R$. For almost every $s\in(r,R)$,
\[
 \partial(T\llcorner B_s(a))
 =(\partial T)\llcorner B_s(a)+\langle T,|x-a|,s\rangle.
\]
There is such an $s$ for which
\[
 \mathbf M\bigl(\langle T,|x-a|,s\rangle\bigr)
 \leq\frac{\|T\|(B_R(a)\setminus\overline B_r(a))}{R-r}.
\]
Consequently the boundary mass of the restricted current is at most this quantity plus $\|\partial T\|(B_R(a))$.
\end{proposition}
\begin{proof}
The distance function is one-Lipschitz. Apply the slicing theorem over the annulus and integrate over $s\in(r,R)$. The integral of the slice masses is no larger than the numerator, so some good level has mass no larger than their average. The restriction-boundary formula is the definition of the slice, and the boundary-mass estimate follows by the triangle inequality. A good level is chosen from the full-measure set where all slicing conclusions hold.
\end{proof}

\begin{example}[The slice of an oriented rectangle]\label{ex:rectangle-slice}
Let $T$ be integration over $[0,1]^2$ with orientation $e_1\wedge e_2$, and set $f(x,y)=x$. For $0<t<1$, $\langle T,f,t\rangle$ is the vertical segment $\{t\}\times[0,1]$ oriented toward increasing $y$. Indeed $(e_1\wedge e_2)\mathbin{\lrcorner}e_1=e_2$. Its boundary is $\delta_{(t,1)}-\delta_{(t,0)}$. The slice of the oriented boundary of the rectangle has the opposite sign, directly illustrating Proposition \ref{prop:slice-boundary-sign}.
\end{example}

\begin{definition}[The integral Plateau problem in a compact container]\label{def:integral-plateau}
Fix $1\leq m<N$, a compact set $K\subset\mathbb R^N$, and an integral $(m-1)$-current $Z$ of finite mass. The admissible class consists of integral $m$-currents $T$ with finite mass, $\operatorname{spt}T\subset K$, and $\partial T=Z$. The Plateau problem minimizes $\mathbf M(T)$ in this class. We assume the class is nonempty when asserting existence; in particular $\partial Z=0$ is necessary. No parametrization, connectedness, or genus is fixed by this definition. Compare the compactness approach in \cite{Sim18}, Chapter 6, Sections 3 and 7.
\end{definition}

\begin{theorem}[Existence of a mass-minimizing integral filling]\label{thm:plateau-direct-method}
If the class in Definition \ref{def:integral-plateau} is nonempty, its mass infimum is attained by an integral current supported in $K$.
\end{theorem}
\begin{proof}
Choose one admissible current to bound the infimum from above, then take a minimizing sequence with bounded masses. Its boundary is always $Z$, so the boundary masses are uniformly bounded. Federer--Fleming compactness, Theorem \ref{thm:integral-compactness}, gives a weakly convergent subsequence with an integral limit supported in $K$. Boundary continuity gives $\partial T=Z$. Mass lower semicontinuity shows that its mass is no larger than the infimum, while admissibility gives the opposite inequality. Thus it attains the infimum. This is the direct-method consequence of the compactness theorem in \cite{Sim18}, Section 3; regularity of the minimizer is a separate problem.
\end{proof}

\begin{corollary}[Compact boundaries have Euclidean minimizers]\label{cor:euclidean-plateau}
Let $1\leq m<N$, and let $Z$ be a compactly supported integral $(m-1)$-current of finite mass. If $m\geq2$, assume $\partial Z=0$; if $m=1$, assume the total signed multiplicity of $Z$ is zero. Then there is a compactly supported integral $m$-current minimizing mass among all compactly supported integral fillings of $Z$ in $\mathbb R^N$.
\end{corollary}
\begin{proof}
Choose a closed ball containing the support of $Z$. Proposition \ref{prop:current-cone-bound}, or Example \ref{ex:zero-current-augmentation} when $m=1$, supplies a filling in that ball. The nearest-point projection $p$ to a closed Euclidean ball is one-Lipschitz and fixes $Z$. Every compactly supported integral filling $S$ therefore projects to another filling $p_\#S$ in the ball, with no larger mass, by Theorem \ref{thm:lipschitz-current-pushforward}. Thus the infimum in the ball is the full compactly supported infimum. Apply Theorem \ref{thm:plateau-direct-method} in this fixed container. The use of compact supports makes properness automatic for the projection.
\end{proof}
